\section{Maximum principles} \label{maxprin}

The maximum-principle statements below are referenced by \cite{Chow-Lu-Ni}.

\begin{theorem}[maximum principles theorem]
  \label{thm:kl-maximum-principles-maximum-principles-theorem}
  \source{kleiner-lott}
  \dcref{Theorem A.1}
  \group{Kleiner--Lott Analytic Appendices}

Let $M$ be a closed manifold, let $\{g(t)\}_{t \in [0,T]}$ be a smooth
one-parameter family of Riemannian metrics, and let $\{X(t)\}_{t \in [0,T]}$
be a smooth one-parameter family of vector fields. Let
$F \: : \: \R \times [0,T] \rightarrow \R$ be Lipschitz. If $u=u(x,t)$ is
$C^2$ in $x$, $C^1$ in $t$, and
\begin{equation}
\frac{\partial u}{\partial t} \: \le \: \triangle_{g(t)}u \: + \: X(t)u \: + \: F(u,t),
\end{equation}
and $\phi \: : \: [0,T] \rightarrow \R$ solves
$\frac{d\phi}{dt} \: = \: F(\phi(t),t)$ with $\phi(0)=\alpha$, then
$u(\cdot,0)\leq\alpha$ implies $u(\cdot,t)\leq\phi(t)$ for all
$t\in[0,T]$.
\end{theorem}

\begin{theorem}[estimates for curvature]
  \label{thm:kl-maximum-principles-estimates-curvature}
  \source{kleiner-lott}
  \dcref{Theorem A.3}
  \group{Kleiner--Lott Analytic Appendices}

Let $(M,g(\cdot))$ be a complete Ricci flow solution on $[0,T]$ with
uniformly bounded curvature. If $u=u(x,t)\in W^{1,2}_{loc}$ weakly satisfies
$\frac{\partial u}{\partial t}\leq\triangle_{g(t)}u$, has
$u(\cdot,0)\leq0$, and, for some $c>0$,
\begin{equation}
\int_0^T\int_M e^{-\:c\:d_t^2(x,x_0)}u^2(x,s)\:dV(x)\:ds<\infty,
\end{equation}
then $u(\cdot,t)\leq0$ for all $t\in[0,T]$.
\end{theorem}

\begin{theorem}[existence for maximum principles]
  \label{thm:kl-maximum-principles-existence-maximum-principles}
  \source{kleiner-lott}
  \dcref{Theorem A.5}
  \group{Kleiner--Lott Analytic Appendices}

Let $M$ be connected, let $\{g(t)\}_{t\in[0,T]}$ be a smooth one-parameter
family of Riemannian metrics, let $\{X(t)\}_{t\in[0,T]}$ be a smooth
one-parameter family of vector fields, and let
$F\: : \:\R\times[0,T]\rightarrow\R$ be Lipschitz. Suppose that $u=u(x,t)$
is $C^2$ in $x$, $C^1$ in $t$, and
\begin{equation}
\frac{\partial u}{\partial t}\: \le \:\triangle_{g(t)}u\:+\:X(t)u\:+\:F(u,t).
\end{equation}
If $\phi\: : \:[0,T]\rightarrow\R$ solves
$\frac{d\phi}{dt}\:=\:F(\phi(t),t)$ and
$u(\cdot,t)\leq\phi(t)$ for all $t\in[0,T]$, while
$u(x_0,t_0)<\phi(t)$ for some $x_0\in M$ and $t_0\in(0,T]$, then there is
$\epsilon>0$ such that $u(\cdot,t)<\phi(t)$ for
$t\in(t_0-\epsilon,T]$.
No completeness assumption on the metrics is required.
\end{theorem}

\begin{theorem}[existence for curvature operators]
  \label{thm:kl-maximum-principles-existence-curvature-operators}
  \source{kleiner-lott}
  \dcref{Theorem A.7}
  \group{Kleiner--Lott Analytic Appendices}
 \label{splitting}

Let $M$ be connected and let $\{g(t)\}_{t\in[0,T]}$ be a Ricci flow with
nonnegative curvature operator. For each $t\in(0,T]$,
$\Image(\Rm_{g(t)})$ is a smooth subbundle of $\Lambda^2(T^*M)$ invariant
under spatial parallel translation. There are times
$0=t_0<t_1<\ldots<t_k=T$ such that, for $1\leq i\leq k$,
$\Image(\Rm_{g(t)})$ is a Lie subalgebra of
$\Lambda^2(T^*_mM)\cong o(n)$ independent of $t\in(t_{i-1},t_i]$, and
$\Image(\Rm_{g(t_i)})\subset\Image(\Rm_{g(t_{i+1})})$.
In particular, a local isometric splitting at one time implies a local
isometric splitting at earlier times.
\end{theorem}

\section{$\phi$-almost nonnegative curvature}
\label{phiappendix}

In dimension three, the Ricci flow equation gives
\begin{equation} \label{3dR}
\frac{dR}{dt} \: = \: \triangle R \: + \: \frac{2}{3}R^2
+2\:|R_{ij}-\frac{R}{3}g_{ij}|^2.
\end{equation}
For a Ricci flow on $[0,T)$ with complete time slices and curvature bounded
on compact time intervals, the maximum principle yields
\begin{equation} \label{Rlowerbound}
(\inf R)(t)\: \ge \:
\frac{(\inf R)(0)}{1-\frac23t(\inf R)(0)}.
\end{equation}
Thus $tR(\cdot,t)>-\frac32$ for $t\geq0$ \cite[Section 2]{Hamiltonn}.

The curvature operator acts on $2$-forms; constant sectional curvature $k$
means multiplication by $2k$, and its trace is the scalar curvature. In
dimension three, nonnegative curvature operator is equivalent to nonnegative
sectional curvature, and each curvature-operator eigenvalue is twice a
sectional curvature.

Assume that at $t=0$ the ordered curvature-operator eigenvalues satisfy
$\lambda_1\leq\lambda_2\leq\lambda_3$ and $\lambda_1\geq-1$. Then
$R(\cdot,t)\geq-\frac32\frac1{t+\frac12}$. For $X=-\lambda_1>0$ at $(x,t)$,
Hamilton--Ivey pinching gives \cite[Theorem 4.1]{Hamiltonn}
\begin{equation}
R(x,t)\: \ge \:X\left(\ln X+\ln(1+t)-3\right),
\end{equation}
or equivalently
\begin{equation} \label{estimate}
tR(x,t)\: \ge \:tX\left(\ln(tX)+\ln(\frac{1+t}{t})-3\right).
\end{equation}

\begin{definition}[positivity for curvature operators]
  \label{def:kl-almost-nonnegative-curvature-positivity-curvature-operators}
  \source{kleiner-lott}
  \lean{KleinerLott.IsHamiltonIveyPinched}
  \dcref{Definition B.5}
  \group{Kleiner--Lott Analytic Appendices}
 \label{HamIvey}

Given $t\geq0$, a Riemannian $3$-manifold $(M,g)$ satisfies the
\emph{time-$t$ Hamilton-Ivey pinching condition} if, for every $x\in M$,
the ordered curvature-operator eigenvalues
$\la_1\leq\la_2\leq\la_3$ satisfy either $\la_1\geq0$, or, with
$\la_1<0$ and $X=-\lambda_1$,
\[tR(x)\geq tX\left(\ln(tX)+\ln\left(\frac{1+t}{t}\right)-3\right).\]
\end{definition}

\begin{lemma}[monotonicity for scalar curvature]
  \label{lem:kl-almost-nonnegative-curvature-monotonicity-scalar-curvature}
  \source{kleiner-lott}
  \dcref{Lemma B.6}
  \group{Kleiner--Lott Analytic Appendices}

\label{lempinchingmonotonic}
If $\Rm$ and $\Rm'$ are $3$-dimensional curvature operators with
$R'\geq R\geq0$ and $\la_1'\geq\la_1$, then the time-$t$ Hamilton-Ivey
pinching condition for $\Rm$ implies that for $\Rm'$.
\end{lemma}
\begin{proof}
\sketch
Set $X=-\la_1$ and $X'=-\la_1'$. If $\la_1'<0$ and
$\ln(tX')+\ln((1+t)/t)-3>0$, the function
\begin{equation}
Y\mapsto tY\left(\ln(tY)+\ln\left(\frac{1+t}{t}\right)-3\right)
\end{equation}
is increasing on the relevant interval, and
\begin{equation}
tX\left(\ln(tX)+\ln\left(\frac{1+t}{t}\right)-3\right)
\geq tX'\left(\ln(tX')+\ln\left(\frac{1+t}{t}\right)-3\right).
\end{equation}
If $\la_1'\geq0$ or the logarithmic factor is nonpositive, the condition is
satisfied directly (in the latter case because $R'>0$).
\end{proof}

For every $s\in\R$, the pinching estimate gives a lower bound
$t\Rm(\cdot,t)\geq\const(s,t)$ whenever $tR(\cdot,t)\leq s$; this is
vacuous for $s\leq-\frac32$. There is a positive
$\Phi\in C^\infty(\R)$ such that $\Phi$ is nondecreasing,
$\Phi(s)/s$ is decreasing for $s>0$, $\Phi(s)\sim s/\ln s$ for large $s$,
and
\begin{equation} \label{phi}
\Rm(\cdot,t)\: \ge \:-\Phi(R(\cdot,t)).
\end{equation}
For an unscaled flow, the precise bound depends on $t_0$ and the time-zero
metric through its lower curvature bound.

\section{Ricci solitons} \label{Riccisolitons}

For a time-dependent family of vector fields $V(t)$, assuming integrability
on a noncompact $M$, the solution of
\begin{equation}
\frac{dg}{dt}\:=\:{\mathcal L}_{V(t)}g
\end{equation}
is
\begin{equation}
g(t)\:=\:\phi^{-1}(t)^*g(t_0),
\end{equation}
where $\phi(t)$ is generated by $-V$ and $\phi(t_0)=\Id$. The convention is
\begin{equation}
{\mathcal L}_{W(t)}\:=\:\eta^{-1}(t)^*\:
\frac{d\eta(t+\epsilon)^*}{d\epsilon}\Big|_{\epsilon=0}
\:=\:-\frac{d\eta^{-1}(t+\epsilon)^*}{d\epsilon}\Big|_{\epsilon=0}\:\eta(t)^*.
\end{equation}

A steady soliton has time-independent $V$ and
\begin{equation}
2\Ric\:+\:{\mathcal L}_Vg\:=\:0,
\end{equation}
with Ricci flow
\begin{equation}
g(t)\:=\:\phi^{-1}(t)^*g(t_0),
\end{equation}
where $\phi(t)$ is generated by $-V$.

A gradient steady soliton satisfies
\begin{align} \label{steady}
\frac{\partial g_{ij}}{\partial t}\: & = \:-2R_{ij}\:=\:
2\nabla_i\nabla_jf,\\
\frac{\partial f}{\partial t}\: & = \:|\nabla f|^2. \notag
\end{align}
Moreover,
\begin{equation}
\frac{\partial}{\partial t}(g^{ij}\partial_jf)\:=\:
2R^{ij}\nabla_jf+g^{ij}\nabla_j|\nabla f|^2
\:=\:-2(\nabla^i\nabla^jf)\nabla_jf+\nabla^i|\nabla f|^2\:=\:0,
\end{equation}
so $V=\nabla f$ is constant in $t$, and
\begin{align} \label{steadysoln}
g(t)\: & = \:\phi^{-1}(t)^*g(t_0),\\
f(t)\: & = \:\phi^{-1}(t)^*f(t_0). \notag
\end{align}
Conversely, if
\begin{equation}
\widehat R_{ij}+\widehat\nabla_i\widehat\nabla_j\widehat f\:=\:0,
\end{equation}
and $V=\widehat\nabla\widehat f$, then
\begin{align}
g(t)\: & = \:\phi^{-1}(t)^*\widehat g,\\
f(t)\: & = \:\phi^{-1}(t)^*\widehat f \notag
\end{align}
satisfy (\ref{steady}). A solution also satisfies
\begin{equation}
\frac{\partial f}{\partial t}\:=\:|\nabla f|^2-\triangle f-R,
\end{equation}
and
\begin{equation}
\frac{\partial}{\partial t}e^{-f}\:=\:-\triangle e^{-f}+Re^{-f}.
\end{equation}

A shrinking soliton is defined on $(-\infty,0)$ and satisfies
\begin{equation}
2\Ric+{\mathcal L}_Vg+\frac gt\:=\:0,
\end{equation}
with $V(t)=-\frac1tV(-1)$ and
\begin{equation} \label{shrinkingflow}
g(t)\:=\:-t\:\phi^{-1}(t)^*g(-1),
\end{equation}
where $\phi(-1)=\Id$ and $\phi$ is generated by $-V$.
A gradient shrinking soliton satisfies
\begin{align} \label{shrinking}
\frac{\partial g_{ij}}{\partial t}\: & = \:-2R_{ij}\:=\:
2\nabla_i\nabla_jf+\frac{g_{ij}}t,\\
\frac{\partial f}{\partial t}\: & = \:|\nabla f|^2. \notag
\end{align}
Here $V=\nabla f$ satisfies $V(t)=-\frac1tV(-1)$, and
\begin{align} \label{shrinkingsoln}
g(t)\: & = \:-t\:\phi^{-1}(t)^*g(-1),\\
f(t)\: & = \:\phi^{-1}(t)^*f(-1). \notag
\end{align}
Conversely, if
\begin{equation}
\widehat R_{ij}+\widehat\nabla_i\widehat\nabla_j\widehat f
-\frac12\widehat g\:=\:0,
\end{equation}
then with $V(t)=-\frac1t\widehat\nabla\widehat f$ the definitions
\begin{align}
g(t)\: & = \:-t\:\phi^{-1}(t)^*\widehat g,\\
f(t)\: & = \:\phi^{-1}(t)^*\widehat f \notag
\end{align}
satisfy (\ref{shrinking}).

An expanding soliton is defined on $(0,\infty)$ and satisfies
\begin{equation}
2\Ric+{\mathcal L}_Vg+\frac gt\:=\:0,
\end{equation}
with $V(t)=\frac1tV(1)$ and
\begin{equation}
g(t)\:=\:t\:\phi^{-1}(t)^*g(1),
\end{equation}
where $\phi(1)=\Id$ and $\phi$ is generated by $-V$.
A gradient expanding soliton satisfies
\begin{align} \label{expanding}
\frac{\partial g_{ij}}{\partial t}\: & = \:-2R_{ij}\:=\:
2\nabla_i\nabla_jf+\frac{g_{ij}}t,\\
\frac{\partial f}{\partial t}\: & = \:|\nabla f|^2. \notag
\end{align}
Here $V=\nabla f$ satisfies $V(t)=\frac1tV(1)$, and
\begin{align} \label{expandingsoln}
g(t)\: & = \:t\:\phi^{-1}(t)^*g(1),\\
f(t)\: & = \:\phi^{-1}(t)^*f(1). \notag
\end{align}
Conversely, if
\begin{equation}
\widehat R_{ij}+\widehat\nabla_i\widehat\nabla_j\widehat f
+\frac12\widehat g\:=\:0,
\end{equation}
then with $V(t)=\frac1t\widehat\nabla\widehat f$ the definitions
\begin{align}
g(t)\: & = \:t\:\phi^{-1}(t)^*\widehat g,\\
f(t)\: & = \:\phi^{-1}(t)^*\widehat f \notag
\end{align}
satisfy (\ref{expanding}). Einstein metrics with $V=0$ give solitons, and
every steady or expanding soliton on a closed manifold is Einstein
\cite[Chapter 2]{Chow-Knopf}.

\section{Local derivative estimates} \label{applocalder}

\begin{theorem}[existence for local derivative estimates]
  \label{thm:kl-local-derivative-estimates-existence-local-derivative-estimates}
  \source{kleiner-lott}
  \dcref{Theorem D.1}
  \group{Kleiner--Lott Analytic Appendices}
 \label{localder}

For $\alpha,K,K',l\geq0$ and $m,n\in\Z^+$, there is
$C=C(\alpha,K,K',l,m,n)$ such that, for $r>0$, if $g(t)$ is a Ricci flow
on $[0,\overline t]$, $0<\overline t\leq\frac{\alpha r^2}{K}$, defined on
an open neighborhood $U$ of $p\in M^n$, with
$\overline{B(p,r,0)}\subset U$ compact and
\begin{equation}
|\Rm(x,t)|\: \le \:\frac K{r^2}
\end{equation}
for $x\in U$, $t\in[0,\overline t]$, and
\begin{equation}
|\nabla^\beta\Rm(x,0)|\: \le \:\frac{K'}{r^{|\beta|+2}}
\end{equation}
for $x\in U$, $|\beta|\leq l$, then
\begin{equation}
|\nabla^\beta\Rm(x,t)|\: \le \:
\frac{C}{r^{|\beta|+2}\left(\frac{t}{r^2}\right)^{\frac{\max(m-l,0)}2}}
\end{equation}
for $x\in B(p,r/2,0)$, $t\in(0,\overline t]$, and $|\beta|\leq m$.
In particular, $|\nabla^\beta\Rm(x,t)|\leq C/r^{|\beta|+2}$ when
$|\beta|\leq l$ \cite{Shi2}; the extension to $l\geq0$ appears in
\cite[Appendix B]{Lu-Tian}.
\end{theorem}

\section{Convergent subsequences of Ricci flow solutions}
\label{subsequence}

\begin{theorem}[existence for convergent subsequences of ricci flow solutions]
  \label{thm:kl-convergent-subsequences-ricci-flow-solutions-existence-convergent-subsequences-ricci-flow-solutions}
  \source{kleiner-lott}
  \dcref{Theorem E.1}
  \group{Kleiner--Lott Analytic Appendices}
\label{accumulate}

Given $r_0\in(0,\infty]$, let $\{g_i(t)\}_{i=1}^\infty$ be Ricci flow
solutions on connected pointed $(M_i,m_i)$ for
$t\in(A,B)$, where $-\infty\leq A<0<B\leq\infty$. Assume
$M_i=B_0(m_i,r_0)$ and $\overline{B_0(m_i,r)}$ is compact for every
$r\in(0,r_0)$. Suppose:
\begin{enumerate}
\item for each $r\in(0,r_0)$ and compact $I\subset(A,B)$, some
$N_{r,I}<\infty$ satisfies
$\sup_{B_0(m_i,r)\times I}|\Rm(g_i)|\leq N_{r,I}$ for all $i$;
\item $\{\inj(g_i(0))(m_i)\}_{i=1}^\infty$ is uniformly bounded below by a
positive number.
\end{enumerate}
After passing to a subsequence, $g_i(t)$ converges smoothly to a Ricci flow
$g_\infty(t)$ on connected pointed $(M_\infty,m_\infty)$ for
$t\in(A,B)$, with $M_\infty=B_0(m_\infty,r_0)$ and compact
$\overline{B_0(m_\infty,r)}$ for $r\in(0,r_0)$. For compact
$I\subset(A,B)$ and $r<r_0$, there are pointed time-independent
diffeomorphisms $\phi_{r,i}:B_0(m_\infty,r)\rightarrow B_0(m_i,r)$ such that
$(\phi_{r,i}\times\Id)^*g_i$ converges smoothly to $g_\infty$ on
$B_0(m_\infty,r)\times I$.
\end{theorem}

Under the sectional-curvature bounds, the injectivity-radius hypothesis is
equivalent to a lower bound on volumes of balls around $m_i$
\cite[Theorem 4.7]{Cheeger-Gromov-Taylor}. The case $r_0=\infty$ and
$N_{r,I}$ independent of $r$ is \cite[Main Theorem]{Hamiltonnn}, as in
Corollary \ref{accumulate3}.

Theorem \ref{accumulate} also applies on $(A,B]$ with
$-\infty\leq A<0\leq B<\infty$. On $[A,B)$, with
$-\infty<A<0<B\leq\infty$, uniform bounds
\begin{equation}
\sup_{B_A(m_i,r)}|\nabla^j\Rm(g_i(A))|\: \le \:C_{r,j}
\end{equation}
give smooth convergence on $[A,B)$ by Appendix \ref{applocalder}; without
these bounds the convergence is $C^0$ on $[A,B)$ and smooth on $(A,B)$.
The analogous statement holds on $[A,B]$, for annuli, and for
$r$-dependent time intervals.

\begin{corollary}[existence for convergent subsequences of ricci flow solutions]
  \label{cor:kl-convergent-subsequences-ricci-flow-solutions-existence-convergent-subsequences-ricci-flow-solutions}
  \source{kleiner-lott}
  \dcref{Corollary E.2}
  \group{Kleiner--Lott Analytic Appendices}
\label{accumulate2}

Let $\{g_i(t)\}_{i=1}^\infty$ be Ricci flows on connected pointed
$(M_i,m_i)$ for $t\in(A,0]$, $-\infty\leq A<0$, with complete time-zero
slices. Suppose:
\begin{enumerate}
\item for every $r\in(0,\infty)$ and compact $I\subset(A,0]$, some
$N_{r,I}<\infty$ satisfies
$\sup_{B_0(m_i,r)\times I}|\Rm(g_i)|\leq N_{r,I}$ for all $i$;
\item for every compact $I\subset(A,0]$, there is $N'_I>0$ such that, for
each $r>0$, some $J_{r,I}\in\Z^+$ satisfies
$\Ric(g_i)\geq-N'_Ig_i$ on $B_0(m_i,r)\times I$ whenever $i\geq J_{r,I}$;
\item $\{\inj(g_i(0))(m_i)\}_{i=1}^\infty$ is uniformly bounded below by a
positive number.
\end{enumerate}
Then a subsequence converges smoothly to a Ricci flow on connected pointed
$(M_\infty,m_\infty)$ for $t\in(A,0]$, with complete time slices.
\end{corollary}
\begin{proof}
\uses{rem:kl-length-distortion-estimates-structural-remark-length-distortion-estimates,thm:kl-convergent-subsequences-ricci-flow-solutions-existence-convergent-subsequences-ricci-flow-solutions}
\sketch
Construct $(M_\infty,m_\infty,g_\infty(\cdot))$ by Theorem
\ref{accumulate}. On each compact $I\subset(A,0]$,
$\Ric(g_\infty)\geq-N'_Ig_\infty$, and by (\ref{similar}),
\begin{equation} \label{similar2}
\dist_t(m_0,m_1)\geq e^{-N'_It}\dist_0(m_0,m_1)
\end{equation}
for $t\in I$. A $g_\infty(t)$-Cauchy sequence is therefore
$g_\infty(0)$-Cauchy and has a $g_\infty(0)$-convergent subsequence. The
identity map is biLipschitz on a sufficiently small $g_\infty(0)$-ball around
its limit, so the original sequence converges in $g_\infty(t)$.
\end{proof}

The forward-time analogue uses flows on $[0,B)$ and, on each compact
$I\subset[0,B)$, eventual bounds $\Ric(g_i)\leq N'_Ig_i$ on
$B_0(m_i,r)\times I$. Under two-sided curvature bounds:

\begin{corollary}[existence for convergent subsequences of ricci flow solutions]
  \label{cor:kl-convergent-subsequences-ricci-flow-solutions-existence-convergent-subsequences-ricci-flow-solutions-3}
  \source{kleiner-lott}
  \dcref{Corollary E.4}
  \group{Kleiner--Lott Analytic Appendices}
\label{accumulate3}

Let $\{g_i(t)\}_{i=1}^\infty$ be Ricci flows on connected pointed
$(M_i,m_i)$ for $t\in(A,B)$, $-\infty\leq A<0<B\leq\infty$, with complete
time-zero slices. Suppose:
\begin{enumerate}
\item for every compact $I\subset(A,B)$ there is $N_I<\infty$ such that,
for each $r>0$, some $J_{r,I}\in\Z^+$ satisfies
$\sup_{B_0(m_i,r)\times I}|\Rm(g_i)|\leq N_I$ whenever $i\geq J_{r,I}$;
\item $\{\inj(g_i(0))(m_i)\}_{i=1}^\infty$ is uniformly bounded below by a
positive number.
\end{enumerate}
Then a subsequence converges smoothly to a Ricci flow on connected pointed
$(M_\infty,m_\infty)$ for $t\in(A,B)$, with complete time slices.
\end{corollary}

When $J_{r,I}=1$ for all $r,I$, this is
\cite[Main Theorem]{Hamiltonnn}.

\section{Harnack inequalities for Ricci flow} \label{appharnack}

Define
\begin{align}
P_{abc}\: & = \:\nabla_aR_{bc}-\nabla_bR_{ac},\\
M_{ab}\: & = \:\triangle R_{ab}-\frac12\nabla_a\nabla_bR
+2R_{acbd}R_{cd}-R_{ac}R_{bc}+\frac{R_{ab}}{2t}. \notag
\end{align}
For a $2$-form $U$ and a $1$-form $W$, set
\begin{equation}
Z(U,W)\:=\:M_{ab}W_aW_b+2P_{abc}U_{ab}W_c+R_{abcd}U_{ab}U_{cd}.
\end{equation}
For a complete Ricci flow with nonnegative curvature operator and curvature
bounded on compact time intervals, $Z(U,W)\geq0$ for $t>0$ and all $U,W$
\cite[Theorem 14.1]{Hamilton}.

With $W_a=Y_a$, $U_{ab}=(X_aY_b-Y_aX_b)/2$, and
\begin{equation} \label{Rict}
\Ric_t(Y,Y)\:=\:(\triangle R_{ab})Y^aY^b+2R_{acbd}R_{cd}Y^aY^b
-2R_{ac}R_{bc}Y^aY^b,
\end{equation}
we have $2Z(U,W)=H(X,Y)$, where
\begin{align} \label{Hexpression}
H(X,Y)\: = & -\Hess_R(Y,Y)-2\langle R(Y,X)Y,X\rangle
+4\left(\nabla_X\Ric(Y,Y)-\nabla_Y\Ric(Y,X)\right)\\
&+2\Ric_t(Y,Y)+2|\Ric(Y,\cdot)|^2+\frac1t\Ric(Y,Y). \notag
\end{align}
For an orthonormal basis $\{e_i\}_{i=1}^n$,
\begin{align}
\sum_iH(X,e_i)\: = & -\triangle R+2\Ric(X,X)
+4\left(\langle\nabla R,X\rangle-\sum_i\nabla_{e_i}\Ric(e_i,X)\right)\\
&+2\sum_i\Ric_t(e_i,e_i)+2|\Ric|^2+\frac1tR. \notag
\end{align}
The traced second Bianchi identity and (\ref{Rict}) give
\begin{equation}
\sum_i\nabla_{e_i}\Ric(e_i,X)\:=\:\frac12\langle\nabla R,X\rangle,
\end{equation}
and
\begin{equation}
\sum_i\Ric_t(e_i,e_i)\:=\:\triangle R.
\end{equation}
Consequently,
\begin{equation} \label{together}
\sum_iH(X,e_i)\:=\:H(X),
\end{equation}
where
\begin{equation} \label{Heqn}
H(X)\:=\:R_t+\frac1tR+2\langle\nabla R,X\rangle+2\Ric(X,X).
\end{equation}
Thus $H(X)\geq0$ for all $X$.

For an ancient solution defined on $(-\infty,0)$, time translation gives
\begin{equation}
R_t+\frac1{t-t_0}R+2\langle\nabla R,X\rangle+2\Ric(X,X)\geq0
\end{equation}
for $t_0\leq t$, and hence
\begin{equation}
R_t+2\langle\nabla R,X\rangle+2\Ric(X,X)\geq0.
\end{equation}
For complete time slices with curvature bounded on compact time intervals,
$R$ is nondecreasing in $t$ and
\begin{equation}
0\leq R_t+2\langle\nabla R,X\rangle+2\Ric(X,X)
\leq R_t+2\langle\nabla R,X\rangle+2R\langle X,X\rangle.
\end{equation}
If $\gamma:[t_1,t_2]\rightarrow M$ is parametrized by $s$ and
$X=\frac12\frac{d\gamma}{ds}$, then
\begin{equation}
\frac{dR(\gamma(s),s)}{ds}=R_t(\gamma(s),s)
+\left\langle\frac{d\gamma}{ds},\nabla R\right\rangle
\geq-\frac12R\left\langle\frac{d\gamma}{ds},\frac{d\gamma}{ds}\right\rangle.
\end{equation}
Therefore, for $t_1<t_2$ and $x_1,x_2\in M$,
\begin{equation}
\label{harnackinequality}
R(x_2,t_2)\geq
\exp\left(-\frac{d_{t_1}^2(x_1,x_2)}{2(t_2-t_1)}\right)R(x_1,t_1).
\end{equation}
For $n=2$, the denominator in the exponential may be replaced by
$4(t_2-t_1)$; if $R(x_2,t_2)=0$, then $g(t)$ is flat for all $t$.

\section{Alexandrov spaces} \label{Alexandrov}

For points $p,x,y$ in a nonnegatively curved Alexandrov space,
$\cangle_p(x,y)$ denotes the angle at $p$ in the Euclidean comparison
triangle. A proper nonnegatively curved Alexandrov space containing a line
splits isometrically as $X=\R\times Y$, with $Y$ proper and nonnegatively
curved \cite[Chapter 10]{Burago-Burago-Ivanov},
\cite{Burago-Gromov-Perelman}. The splitting conclusion is
\cite[Theorem 10.5.1]{Burago-Burago-Ivanov}.

If $M$ is an $n$-dimensional nonnegatively curved Alexandrov space, $p\in M$,
and $\la_k\ra0$, then $(\la_kM,p)$ converges in the pointed
Gromov--Hausdorff sense to the Tits cone $\ctits M$, the Euclidean cone over
$\tits M$, which is nonnegatively curved and has dimension at most $n$
\cite[p. 58-59]{Ballmann-Gromov-Schroeder}. If $\ctits M$ splits as
$\R^k\times Y$, then triangle comparison produces $k$ orthogonal lines
through a basepoint, and the splitting theorem gives an $\R^k$ factor of
$M$.

If $d(x_k,p)\ra\infty$ and $r_k>0$ satisfies
$r_k/d(x_k,p)\ra0$, then $(\frac1{r_k}M,x_k)$ subconverges to a pointed
Alexandrov space $(N_\infty,x_\infty)$ splitting as
$\R\times N'$. Indeed, choose $y_k$ with
$d(y_k,x_k)/d(x_k,p)\ra1$ and $\cangle_{x_k}(p,y_k)\ra\pi$. For every
$\rho<\infty$, choose $p_k\in\ol{px_k}$ and $z_k\in\ol{x_ky_k}$ with
$d(x_k,p_k)/r_k\ra\rho$, $d(x_k,z_k)/r_k\ra\rho$; comparison-angle
monotonicity \cite[Chapter 4.3]{Burago-Burago-Ivanov} gives
$\cangle_{x_k}(p_k,z_k)\ra\pi$. The limit contains a line through
$x_\infty$, so the splitting theorem applies.

If $M$ is complete Riemannian with nonnegative sectional curvature and
$C\subset M$ is compact, connected, and has weakly convex boundary, then
$C_t=\{x\in C\mid d(x,\partial C)\geq t\}$ is convex in $C$
\cite[Chapter 8]{Cheeger-Ebin}. If the second fundamental form of
$\partial C$ is at least $1/r$, then $d(x,\partial C)\leq r$ for every
$x\in C$.

If $M$ is an $n$-manifold, $r_2\geq r_1>0$, $B(p,r_2)$ has compact closure,
and its sectional curvature is bounded below by $K$, then
\begin{equation} \label{BG}
\frac{\vol(B(p,r_2))}{\vol(B(p,r_1))}\: \le \:
\begin{cases}
\frac{\int_0^{r_2}\sin^{n-1}(kr)\:dr}{\int_0^{r_1}\sin^{n-1}(kr)\:dr}
& \text{if }K=k^2,\\
\frac{r_2^n}{r_1^n}& \text{if }K=0,\\
\frac{\int_0^{r_2}\sinh^{n-1}(kr)\:dr}{\int_0^{r_1}\sinh^{n-1}(kr)\:dr}
& \text{if }K=-k^2.
\end{cases}
\end{equation}
For $K=k^2$, $k>0$, assume $r_2\leq\pi/k$. The same inequality holds
under the Ricci lower bound $(n-1)K$ and for Alexandrov curvature bounded
below by $K$.

\section{Finding balls with controlled curvature} \label{pointpicking}

\begin{lemma}[existence for curvature]
  \label{lem:kl-finding-balls-controlled-curvature-existence-curvature}
  \source{kleiner-lott}
  \dcref{Lemma H.1}
  \group{Kleiner--Lott Geometric Appendices}
\label{selection}

Let $X$ be Riemannian with $R\geq0$ and compact
$\ol{B(x,5r)}$. There is a ball $B(y,\bar r)\subset B(x,5r)$ with
$\bar r\leq r$ such that $R(z)\leq2R(y)$ for $z\in B(y,\bar r)$ and
$R(y)\bar r^2\geq R(x)r^2$.
\end{lemma}
\begin{proof}
\sketch
Set $(x_1,r_1)=(x,r)$. If $R\leq2R(x_i)$ on $B(x_i,r_i)$, keep
$(x_{i+1},r_{i+1})=(x_i,r_i)$. Otherwise choose
$x_{i+1}\in B(x_i,r_i)$ with $R(x_{i+1})>2R(x_i)$ and set
$r_{i+1}=r_i/\sqrt2$. All balls lie in $B(x,5r)$; compactness makes the
sequences eventually constant, and the limiting $(y,\bar r)$ has the stated
properties.
\end{proof}

The same selection principle has a spacetime version.

\section{Statement of the geometrization conjecture} \label{geom}

Let $M$ be a connected orientable closed (compact boundaryless) $3$-manifold.
The geometrization conjecture asserts that $M$ is a connected sum of closed
$3$-manifolds $\{M_i\}_{i=1}^n$, each containing a codimension-$0$ compact
submanifold-with-boundary $G_i$ such that
\begin{itemize}
\item $G_i$ is a graph manifold;
\item $M_i-G_i$ admits a complete finite-volume metric of constant negative
sectional curvature;
\item every component $T$ of $\partial G_i$ is an incompressible torus in
$M_i$, i.e. for $t\in T$, $\pi_1(T,t)\rightarrow\pi_1(M_i,t)$ is injective.
\end{itemize}
Here $G_i$ may be $\emptyset$ or $M_i$.

A graph manifold is obtained from pairs of pants $\{P_i\}_{i=1}^N$ and
closed disks $\{D_j^2\}_{j=1}^{N'}$ by taking
$\{S^1\times P_i\}\cup\{S^1\times D_j^2\}$, pairing an even number of
toral boundary components by homeomorphisms, and gluing; assume the gluing
produces an orientable manifold. All graph manifolds arise this way. Its
boundary, when nonempty, is a disjoint union of tori, gluing graph manifolds
along boundary tori gives a graph manifold, and a connected sum is a graph
manifold exactly when each factor is one \cite[Chapter 2.4]{Matveev},
\cite[Proposition 2.4.3]{Matveev}.

Equivalently, use the Kneser--Milnor decomposition into uniquely determined
prime factors, each either $S^1\times S^2$ or irreducible, where every
embedded $S^2$ bounds a $3$-ball. An irreducible $M$ has a minimal
isotopy-unique collection of disjoint incompressible embedded tori
$\{T_k\}_{k=1}^K$ such that the metric completion $M'$ of each component of
$M-\bigcup_{k=1}^K T_k$ is either Seifert or non-Seifert with every embedded
incompressible torus isotopic into $\partial M'$ \cite{Scott}. The latter
case is conjectured to have hyperbolic interior; this holds when
$\partial M'\neq\emptyset$. An orientable Seifert $3$-manifold is a graph
manifold \cite[Proposition 2.4.2]{Matveev}, and every graph manifold is a
connected sum of prime graph manifolds which split along incompressible tori
into Seifert manifolds \cite[Proposition 2.4.7]{Matveev}; hence the two
formulations are equivalent \cite{Scott,Thurston}.
